% !PRINTED 253
\weberpage{253}
The group $E$ contains a divisor of index $p^n + 1$, which consists of all substitutions of the form
\[ \begin{pmatrix} \alpha, & \beta \\ 0, & \delta \end{pmatrix} \]
in which $\beta$ denotes an arbitrary number in $\mathfrak{C}$, and $\alpha$ any number different from zero, and $\delta = \alpha^{-1}$.

To this divisor there corresponds a permutation group on $p^n + 1$ letters isomorphic to the group $E$, which one can form as follows:

The linear expression
\weq{12}{\eta = \dfrac{\alpha\xi + \beta}{\gamma\xi + \delta} \quad \text{or} \quad \xi = \dfrac{\delta\eta - \beta}{-\gamma\eta + \alpha}}
gives for every number $\xi$ in $\mathfrak{C}$ a corresponding number $\eta$, except when $\gamma\xi + \delta = 0$. Likewise, to every $\eta$ there is a definite $\xi$, except for $\gamma\eta - \alpha = 0$. To remove this exception, it suffices to extend the system $\mathfrak{C}$ by adjoining \gesperrt{one} element, which we call ``infinity'' and denote by $\infty$. With this symbol one reckons according to the following rules, in which $\alpha$ denotes any element of the extended system $E$:
% [sic] Weber prints ordinary E both here and in "system E" below, although both references mean the extended field \mathfrak{C}; preserved as printed. Ordinary E otherwise denotes the congruence group.
\[
\begin{array}{l@{\quad}l}
\alpha + \infty = \alpha - \infty = \infty, & \text{except for } \alpha = \infty \\
\alpha \cdot \infty = \infty & \text{\textquotedbl}\ \ \text{\textquotedbl}\ \ \alpha = 0 \\
\alpha : 0 = \infty & \text{\textquotedbl}\ \ \text{\textquotedbl}\ \ \alpha = 0 \\
\alpha : \infty = 0 & \text{\textquotedbl}\ \ \text{\textquotedbl}\ \ \alpha = \infty.
\end{array}
\]
Then the result of every calculation with the four arithmetic operations always leads to a definite number of the system $E$, and only the combinations $\infty \pm \infty, 0 . \infty, 0 : 0, \infty : \infty$ remain undefined.

By (12), to the number $\xi = -\delta : \gamma$ there then corresponds the number $\eta = \infty$, and to the number $\xi = \infty$ the number $\eta = \alpha : \gamma$, and to every substitution $A$ in $E$ there corresponds a definite permutation of the $p^n + 1$ elements of the extended system $\mathfrak{C}$. One sees easily that only the identical substitution leaves all these elements unchanged, and that consequently two different substitutions in $E$ always give rise to different permutations. The isomorphism is therefore one-to-one.
